\documentclass[11pt,a4paper]{article}
\usepackage[utf8]{vietnam}
\usepackage{amsmath,amssymb}
\usepackage{tabularx}
\usepackage{tikz}
\usepackage{xcolor}
\usepackage{fancyhdr}
\usepackage[top=1.8cm, bottom=2.0cm, left=1.8cm, right=1.8cm]{geometry}

\setlength{\parskip}{3pt}
\setlength{\parindent}{0pt}

\pagestyle{fancy}
\fancyhf{}
\renewcommand{\headrulewidth}{0.4pt}
\lhead{\footnotesize \textbf{CHÙA HƯƠNG SƠN -- HƯỚNG DẪN CHẤM THI}}
\rhead{\footnotesize \textbf{ĐÁP ÁN CHI TIẾT ĐỀ THI THỬ SỐ 07}}
\cfoot{\thepage}

\begin{document}

\begin{center}
{\large\bfseries ĐÁP ÁN \& HƯỚNG DẪN CHẤM ĐỀ THI THỬ CUỐI HỌC KỲ I (ĐỀ SỐ 07)}\\[1.5mm]
{\bfseries ĐƠN VỊ: CHÙA HƯƠNG SƠN \quad|\quad MÔN: TOÁN -- LỚP 9 \quad|\quad NĂM HỌC 2026 -- 2027}\\[1mm]
\textit{(Chuẩn cấu trúc định dạng mới của Sở Giáo dục và Đào tạo)}
\end{center}

\vspace{1mm}
\hrule
\vspace{3mm}

\noindent
\textbf{\large PHẦN I. TRẮC NGHIỆM KHÁCH QUAN (7,0 ĐIỂM)}

\vspace{2mm}
\noindent
\textbf{1. Bảng đáp án Phần 1: Trắc nghiệm nhiều lựa chọn (3,0 điểm)} \textit{(Mỗi câu đúng được 0,25 điểm)}
\begin{center}
\begin{tabular}{|c|c|c|c|c|c|c|c|c|c|c|c|c|}
\hline
\textbf{Câu} & \textbf{1} & \textbf{2} & \textbf{3} & \textbf{4} & \textbf{5} & \textbf{6} & \textbf{7} & \textbf{8} & \textbf{9} & \textbf{10} & \textbf{11} & \textbf{12} \\
\hline
\textbf{Đáp án} & \textbf{A} & \textbf{A} & \textbf{A} & \textbf{A} & \textbf{A} & \textbf{A} & \textbf{A} & \textbf{A} & \textbf{A} & \textbf{A} & \textbf{A} & \textbf{A} \\
\hline
\end{tabular}
\end{center}

\vspace{2mm}
\noindent
\textbf{2. Bảng đáp án Phần 2: Trắc nghiệm Đúng -- Sai (2,0 điểm)}
\begin{center}
\begin{tabular}{|c|c|c|c|c|l|}
\hline
\textbf{Câu} & \textbf{Ý a} & \textbf{Ý b} & \textbf{Ý c} & \textbf{Ý d} & \textbf{Ghi chú đáp án chuẩn} \\
\hline
\textbf{Câu 13} & \textbf{Đ} & \textbf{Đ} & \textbf{S} & \textbf{Đ} & Ý c) sai vì sản lượng năm nay là 882 tạ (đề ghi 800 tạ). \\
\hline
\textbf{Câu 14} & \textbf{Đ} & \textbf{Đ} & \textbf{Đ} & \textbf{S} & Ý d) sai: Diện tích hình thang $S = \dfrac{(24+10) \times 7}{2} = 119\text{ cm}^2$ (đề ghi 140). \\
\hline
\end{tabular}
\end{center}

\vspace{2mm}
\noindent
\textbf{3. Bảng đáp án Phần 3: Trắc nghiệm trả lời ngắn (2,0 điểm)} \textit{(Mỗi câu đúng được 0,5 điểm)}
\begin{center}
\begin{tabular}{|c|l|c|c|}
\hline
\textbf{Câu} & \textbf{Nội dung câu hỏi} & \textbf{Kết quả bài làm} & \textbf{Điểm} \\
\hline
\textbf{Câu 15} & Nghiệm bất phương trình $3(x - 4) + 5 > 2(2x + 1) - 15$ & \textbf{$x < 6$} & 0,5đ \\
\hline
\textbf{Câu 16} & Giá trị biểu thức $E = \sqrt{(\sqrt{15}-4)^2} + \sqrt{15}$ & \textbf{$4$} & 0,5đ \\
\hline
\textbf{Câu 17} & Chiều dài dây cáp neo ($100 / \sin 60^\circ$) & \textbf{$115{,}47\text{ m}$} & 0,5đ \\
\hline
\textbf{Câu 18} & Diện tích hình quạt tròn ($\dfrac{3{,}14 \times 10^2 \times 108}{360}$) & \textbf{$94{,}2\text{ cm}^2$} & 0,5đ \\
\hline
\end{tabular}
\end{center}

\vspace{3mm}
\noindent
\textbf{\large PHẦN II. TỰ LUẬN (3,0 ĐIỂM)}

\vspace{2mm}
\textbf{Câu 19 (1,0 điểm).}
\begin{itemize}
  \item[a)] \textit{(0,5 điểm)} $\begin{cases} 3x - 4y = 2 \\ 2x + y = 5 \end{cases} \Leftrightarrow \begin{cases} 3x - 4(5 - 2x) = 2 \\ y = 5 - 2x \end{cases} \Leftrightarrow \begin{cases} 11x = 22 \\ y = 5 - 2x \end{cases} \Leftrightarrow \begin{cases} x = 2 \\ y = 1 \end{cases}$. \hfill \textbf{(0,5đ)}
  \item[b)] \textit{(0,5 điểm)} $A = \dfrac{\sqrt{x}(\sqrt{x}-2) + 2(\sqrt{x}+2)}{x - 4} : \dfrac{x+4}{x-4} = \dfrac{x - 2\sqrt{x} + 2\sqrt{x} + 4}{x - 4} \cdot \dfrac{x - 4}{x + 4} = \dfrac{x + 4}{x + 4} = 1$. \hfill \textbf{(0,5đ)}
\end{itemize}

\vspace{1mm}
\textbf{Câu 20 (1,5 điểm).}
\begin{itemize}
  \item[a)] \textit{(0,75 điểm)} Vì $CA \perp OA, CD \perp OD \Rightarrow \widehat{CAO} = \widehat{CDO} = 90^\circ \Rightarrow A, C, D, O$ cùng thuộc đường tròn đường kính $OC$. Lại có $CA = CD$ và $OA = OD \Rightarrow OC$ là đường trung trực của $AD \Rightarrow OC \perp AD$ tại $I$. \hfill \textbf{(0,75đ)}
  \item[b)] \textit{(0,75 điểm)} Vì $DE$ là đường kính của $(O) \Rightarrow \widehat{DBE} = 90^\circ \Rightarrow EB \perp AD$. Lại có $OC \perp AD \Rightarrow EB \parallel OC$. Tam giác vuông $ACD \backsim \Delta EBD \Rightarrow AC \cdot EB = R^2$ (hoặc dùng hệ thức lượng trong tam giác vuông $ACO$). \hfill \textbf{(0,75đ)}
\end{itemize}

\vspace{1mm}
\textbf{Câu 21 (0,5 điểm).}
\begin{itemize}
  \item[a)] \textit{(0,25 điểm)} Gọi chiều rộng đáy $x > 0 \Rightarrow$ chiều dài đáy $2x$. Chiều cao $h > 0$. $V = 2x^2 h = 500 \Rightarrow h = \dfrac{250}{x^2}$. Diện tích toàn phần:
  $S_{tp} = 2(2x^2) + 2(x \cdot h + 2x \cdot h) = 4x^2 + 6xh = 4x^2 + \dfrac{1500}{x} = 4x^2 + \dfrac{750}{x} + \dfrac{750}{x} \ge 3\sqrt[3]{4x^2 \cdot \dfrac{750}{x} \cdot \dfrac{750}{x}} = 3\sqrt[3]{2250000} \approx 393\text{ cm}^2$. Dấu ``='' khi $4x^2 = \dfrac{750}{x} \Leftrightarrow x^3 = \dfrac{750}{4} = 187{,}5$. \hfill \textbf{(0,25đ)}
  \item[b)] \textit{(0,25 điểm)} Áp dụng BĐT Cauchy-Schwarz dạng Engel:
  \[
  P = \dfrac{1}{a} + \dfrac{1}{b} + \dfrac{1}{c} \ge \dfrac{(1+1+1)^2}{a+b+c} = \dfrac{9}{1} = 9.
  \]
  Dấu ``='' xảy ra khi $a = b = c = \dfrac{1}{3}$. Vậy $\min P = 9$. \hfill \textbf{(0,25đ)}
\end{itemize}

\end{document}
